\section{Introduction}

Generative models are commonly regarded as more challenging than discriminative models.
While discriminative modeling typically focuses on mapping individual samples to their corresponding labels, generative modeling concerns mapping from one distribution to another.
This can be expressed as learning a mapping $f$ such that the \mbox{\textit{pushforward}} of a prior distribution $p_\text{prior}$ matches the data distribution, namely, $f_\# p_\text{prior} \approx p_{\text{data}}$.
Conceptually, generative modeling learns a \textit{functional} (here, $f_\#$) that maps from one function (here, a distribution) to another.

The ``pushforward'' behavior can be realized \textit{iteratively} at \textit{inference} time, \eg, in prevailing paradigms such as Diffusion~\cite{sohl2015deep} and Flow Matching~\cite{lipman2022flow}. When generating, these models map noisier samples to slightly cleaner ones, progressively evolving the sample distribution toward the data distribution.
This modeling philosophy can be viewed as decomposing a complex pushforward map (\ie, $f_\#$) into a chain of more feasible transformations, applied at inference time.

In this paper, we propose \emph{Drifting Models}, a new paradigm for generative
modeling. Drifting Models are characterized by learning a pushforward map that evolves during \textit{training} time, thereby removing the need for an iterative inference procedure. The mapping $f$ is represented by a single-pass, non-iterative network. As the training process is inherently iterative in deep learning optimization, it can be naturally viewed as evolving the pushforward
distribution, $f_\# p_\text{prior}$, through the update of $f$. See Fig.~\ref{fig:teaser}.

To drive the evolution of the training-time pushforward, we introduce a
\emph{drifting field} that governs the sample movement.
This field depends on the generated distribution and the data distribution.
By definition, this field becomes zero when the two distributions match, thereby reaching an equilibrium in which the samples no longer drift.

Building on this formulation, we propose a simple training objective that minimizes the \textit{drift} of the generated samples. This objective induces sample movements and thereby evolves the underlying pushforward distribution through iterative optimization (\eg, SGD). We further introduce the designs of the drifting field, the neural network model, and the training algorithm.

Drifting Models naturally perform \textit{single-step} (\mbox{``1-NFE''}) generation and achieve strong empirical performance.
On ImageNet 256$\times$256, we obtain a 1-NFE
FID of \textbf{1.54} under the standard latent-space generation protocol, achieving a new state-of-the-art among single-step methods. This result remains competitive even when compared with \textit{multi-step} diffusion-/flow-based models.
Further, under the more challenging \textit{pixel}-space generation protocol (\ie, without latents), we reach a 1-NFE FID of \textbf{1.61}, substantially outperforming previous pixel-space methods. 
These results suggest that Drifting Models offer a promising new paradigm for high-quality, efficient generative modeling.
